\chapter{Discussion and Limitations}
\label{ch:discussion}
\section{Answers to the research questions}
\label{sec:rq-answers}

\subsection{RQ1: detection and representation}

The count representation provides useful attack evidence in the selected
traces, with substantial variation across devices. The closing replay's
pooled F1 is 80.00\%, but Chromecast achieves 61.54\% and Samsung
90.91\%. Eight of nine fully observed attack intervals are detected,
and seven receive an alert by attack end. These results support the
feasibility of a compact, device-specific monitor for the selected devices,
attacks, and observation periods.

The retained count tails are useful safeguards against losing alerts
when historical adaptive references move. Static confusion counts remain
unchanged on the inspected support. Consequently, the
joint rule's justification is broader decision coverage within the
same compact feature set, with observed protection in legacy adaptive runs.

Window duration has a direct effect on the task. Coarser windows can
increase duration-specific F1 and interval coverage while delaying
alerts beyond attack end. Window metrics and alert availability therefore
need joint interpretation. Ten minutes is the practical evaluated reference
for this project, with an explicit delay cost.

\subsection{RQ2: numerical computation}

LM and SciPy preserve all tested development decisions across three
datasets, while fitted parameters and scores remain close. The native
method is successfully integrated without changing the row-wise
statistical objective. Its runtime benefit is operation-dependent:
single-window scoring is faster in the retained checks, while the
production matrix kernel and complete fitting favour SciPy. The
aggregate professor-style experiment favours LM on its own workload,
which has a different objective boundary and initialisation.

The supported conclusion is therefore accurate numerical integration
with operation-dependent performance. Attack detection remains governed by
the same probability and decision rule. The distinction between estimator,
likelihood backend, and detector is necessary to interpret the result
and to attribute the inherited work correctly
\cite{minka2000,languasco2023}.

\subsection{RQ3: adaptation and safeguards}

Historical refitting improves Chromecast's development F1 but fails
to improve Samsung's. Large Samsung concentration changes demonstrate
that a refit can move well beyond the original dispersion. The retained
conservative rules reject all closing replay candidates, preserving
the initial model/calibration pair and static decisions. This demonstrates
enforceable rejection; successful benign-drift adaptation remains a future
experimental target.

The duration experiment supplies a further negative result: accepted
conservative threshold updates can increase false alerts severely at
other durations. Safeguards reduce specific permitted movements, while
long-run quality remains an empirical outcome. This result is consistent
with the broader difficulty of separating drift from adversarial
contamination in unsupervised streams \cite{gama2014,kloft2012}.
The thesis establishes an inspectable experimental framework for
that tradeoff, with both accepted and rejected candidates visible.

\section{What the four counts preserve and discard}
\label{sec:representation-limits}

\subsection{Statistical anomalies require security interpretation}

A fitted normal model measures compatibility with its reference. A benign
software update, changed application use, or a new
network service can shift counts. An attack can leave the selected
mixture and count tails inside their normal ranges. Both directions
of error follow from the mismatch between statistical abnormality
and security intent, a central limitation of network anomaly
detection discussed by Sommer and Paxson \cite{sommer2010}.

The device-specific design makes this question more concrete than a
single model of arbitrary network traffic. Normal references do differ
strongly across the three workloads. A single device can still have
several changing benign uses. The initial captures
represent a limited set of operating conditions. The saved profile
can be an intelligible reference without being an exhaustive model
of every future legitimate behaviour.

An event therefore retains evidence for the configured detector decision.
Counts, expected shares, margins, and model identity explain why it alerted.
Exploit attribution, endpoint attribution, and response selection require
additional evidence. Attack types in the result tables come from offline
dataset annotations.

\subsection{Several packet sequences can have exactly the same input}

Aggregation maps many packet sequences to the same four-dimensional
vector. Two windows can contain the same TCP, UDP, SSDP, and ARP counts
while differing in destinations, TCP flags, payloads, packet sizes,
and timing within the interval. The implemented detector must assign
the same decision to those windows when its active state is the
same. Recovering the discarded information requires a richer input.

This explains a specific limitation for TCP SYN attacks. TCP SYN
packets and ordinary TCP packets increment the same category. Their
distinction is visible only indirectly if the attack changes the
aggregate enough to cross a count or composition cutoff. Low-rate
behaviour that resembles the existing reference can remain undetected,
as in the missed Chromecast reflection interval. This is a plausible
representation-based explanation; one interval cannot establish a unique
causal mechanism.

The exclusive SSDP category is helpful because discovery traffic
would otherwise be merged with ordinary UDP. It still discards
message contents, request/response roles, and destination structure.
Similarly, an unusual ARP volume is visible, but malicious ARP
semantics with an ordinary count need not be. Extending these signals
would require a revised input contract, prepared datasets, models,
and evaluation.

\subsection{Conditional probability and volume are different signals}

The DM model gives a probability mass for a composition conditional
on the observed total $N$. Its complete log probability includes the
multinomial coefficient in addition to the relationship between the
normalised observed vector and the fitted mean. The number and
probability of possible compositions vary with $N$. Raw score values
can therefore reflect both composition and the conditional sample
space without defining a separate probability model for $N$.

The total-count guard addresses a different question: whether the
window's total is unusually high relative to benign counts. The ARP
guard adds a particular absolute-category check. Their joint
calibration avoids claiming that three separately chosen one-percent
tests automatically form a one-percent detector. Its future false-alert
rate depends on later traffic and temporal dependence.

The retained guards are upper tails. A silent observed vector does
not violate them, and the zero-count conditional composition has no
internal evidence of a changed mixture. A device going offline is
thus a distinct monitoring question. It could require an availability
signal based on expected activity or missingness, rather than an
interpretation that zero packets is necessarily an attack. Preserving
the difference between observed silence and absent measurement is
essential to any such extension.

\section{Adaptation is a selection problem as well as a fitting problem}
\label{sec:adaptation-discussion}

\subsection{Learning buffers contain statistically admitted observations}

An accepted learning window is one that passes implemented checks.
The term \emph{safe buffer} describes its admission policy; benign status
remains unknown. The original
references and count guards make admission stricter, but an attack
whose counts resemble normal behaviour can still pass. Conversely,
a legitimate new behaviour can fail the same checks and never enter
the buffer needed to adapt toward it.

This creates an unavoidable tension in the present design. Strong
rejection protects the trusted reference but can prevent adjustment
to benign change. Weak rejection increases the opportunity to
adapt but permits contaminated or unrepresentative observations to
change the detector. Work on poisoning and normal-profile construction
motivates explicit attention to this tradeoff
\cite{kloft2012,cretu2008}. The project implements concrete checks and
evaluates their observed effect.

The buffer is also selected even when every accepted observation is
benign. Restricting admission to observations with large normality
margins samples a central subset of the active distribution. Fitting
that subset can underestimate tail variation and increase the
concentration parameter. The dramatic Samsung changes are consistent
with this concern. Separating selection effects from contamination and
actual traffic changes requires a controlled
study would need known benign drift and a separately manipulated
contamination fraction.

\subsection{Candidate checks answer local questions}

The implemented movement limits answer whether a candidate is close
enough to its reference under specified statistics. Later labelled
observations answer whether attacks become easier to detect. A candidate can
increase its own buffer likelihood because the same selected buffer
was used for fitting. That diagnostic checks numerical fit behaviour;
independent validation requires separate future observations.

Threshold movement also needs a time horizon. Bounding each accepted
step can still permit substantial movement after several steps. A
fixed original composition reference and count limits constrain
some ways of losing sensitivity. The duration study's Samsung false
positives show that an adapting DM cutoff can create extra alerts.
A policy described as conservative should therefore be reported
through its equations and acceptance history, not through its name.

Every proposed update in the closing replay is rejected, so this partition
demonstrates preservation of the original reference. Demonstrating useful
adaptation requires controlled benign drift and accepted candidates that
improve later outcomes.

\subsection{Persistence protects state consistency}

Atomic activation ensures that a complete accepted model/calibration
pair becomes current together. Rollback and restart validation prevent
an interrupted write from silently combining incompatible artifacts.
Candidate quality remains governed by statistical validation and later
outcomes.

Likewise, restart resets learning buffers and counters rather than
restoring an exact input-processing position. That choice bounds
state and simplifies recovery but changes the subsequent update
schedule. The external producer must manage replay and duplicate
delivery. Exactly-once processing would require a shared input-position and
acknowledgement protocol.

The local writer lock is sufficient for the tested single-host
process contract. A distributed filesystem or independently managed
replicas require a new ownership and delivery design and their own
measurements.

\section{Interpreting numerical accuracy and runtime}
\label{sec:numerical-discussion}

\subsection{Numerical agreement is necessary but task-specific}

Large log-gamma terms can be individually well represented while
their difference loses significant digits. The numerical method and
fallback address that calculation, following established concerns
about floating-point arithmetic \cite{goldberg1991}. Parameter
validation, native-resource management, and preservation of the
multinomial coefficient are equally necessary for computing the intended
probability on valid inputs.

The LM--SciPy results have two practical implications. First, both backends
produce the same binary decisions on the checked development inputs.
Second, the small observed score differences still matter near a
decision boundary, so the comparison should remain part of future
changes to numerical code. Stress probes extend coverage to large count
vectors and concentrations.

Runtime must also be attributed to the correct layer. Single-window
scoring, repeated matrix evaluation during estimation, and the full
application include different setup, allocation, and iteration costs.
The aggregate experiment changes the likelihood boundary, so its speedup
applies to that workload. Reporting these
layers separately is more useful for implementation decisions than
one averaged ratio that hides their differences.

\subsection{A rejected candidate can dominate latency}

The simulation's mean processing time is small, but its longest
window waits several seconds for a refit that is then rejected.
An operator interested in response time needs this tail measurement
as well as the mean. Rejecting non-convergence protects model
state while leaving the computational cost already incurred.

The synchronous runtime makes that cost visible as backpressure:
it finishes one record before consuming the next. This avoids an
unbounded application queue under the supplied-window contract, but
it shifts responsibility upstream when input arrives faster than
processing. Lossless collection under arbitrary packet rates requires a
collector measurement. An asynchronous implementation
would have different queuing and failure behaviour and would need
its own evaluation.

At the selected ten-minute duration, waiting for observation close
is generally a much larger contribution to attack-alert delay than
the normal scoring computation. Reducing a microsecond-scale kernel
time therefore does little to recover the minutes lost to a coarse
observation. Numerical efficiency and operational timeliness remain
separate dimensions of the system.

\section{Threats to validity and generalisability}
\label{sec:validity}

\subsection{Selection, chronology, and prior label exposure}

Chronological separation prevents future windows from entering the
initial fit. Evaluation bias can still arise when later labels influence
feature, policy, or duration choices. The retained
later partition was inspected in previous audits before the closing
fixed replay. This is the main limitation on interpreting its F1
as a generalisation estimate, in line with the dataset-selection
concerns of Pendlebury et al. and Arp et al.
\cite{pendlebury2019,arp2022}.

The closing protocol improves reproducibility: it records fixed
inputs, settings, predictions, and unchanged starting artifacts. Prior
development knowledge remains part of the study history. An untouched
temporal extension or external collection is required
to assess the selected rule independently. Replaying the same bytes
with more random seeds would not supply that missing independence.

The labelled evaluation contains two device views. Variation across
manufacturers, firmware versions, network topologies, and applications
requires broader deployment-level sampling. The D-Link data remain an
algorithm-benchmark workload. Detection conclusions are therefore expressed
for the observed Chromecast and Samsung partitions.

\subsection{Labels and observation coverage}

Attack labels come from official time intervals, while the detector
classifies aggregates. A window is positive whenever it overlaps an
attack, even if the overlap is brief or contains little selected
traffic. That convention is reproducible but can penalise a normal
aggregate near an interval edge or reward a delayed alert covering
only part of an attack. Interval coverage and first-alert time help
expose these effects and complement the window confusion counts
\cite{tatbul2018}.

Truncated captures are excluded, and partial dates contribute only
observed windows. This avoids inventing counts but changes the
available attack set. The six Chromecast annotations on the excluded capture
remain outside the reported results. The D-Link algorithm benchmark reports
846 retained rows and 702 distinct UTC bounds, with both numerical backends
receiving the identical matrix.

The benign reference is presumed benign from its dataset role.
Normal-only training remains conditional on that trust assumption. A
forensically established reference would require independent verification
of legitimate modes and unannotated activity.

\subsection{Small denominators and dependent observations}

Only 18 windows are positive in the closing replay. Changing a few
predictions can materially change recall and F1. Adjacent windows
and windows from one attack are dependent; treating all 1,008
observations as independent Bernoulli trials would overstate the
amount of independent evidence. Repeated mode outputs are also
paired decisions on the same input, not replication.

For these reasons the report gives exact counts and avoids narrow
confidence statements based on an unsupported independence model.
A larger study could estimate uncertainty using independent devices
or capture blocks and report variability across those units. The
present observations support descriptive comparisons within this study.

Rare positives also explain why a low false positive rate can coexist
with modest precision. Chromecast's three false alerts are few
relative to 426 benign windows, but substantial relative to four
true alerts. Both precision and recall are therefore retained,
rather than using accuracy dominated by true negatives
\cite{saito2015}. Pooling retains the underlying class-imbalance and
prevalence issues.

\subsection{Comparison scope and implementation history}

Related systems establish that learned normal IoT profiles and
online anomaly detection are established research directions
\cite{meidan2018,mirsky2018,nguyen2019}. Direct baselines in this thesis are
the stated static, adaptive, scoring, and numerical alternatives. A ranking
against the neural and ensemble systems would require common data and tuning
budgets.

Source revisions also matter. The closing replay predates the
count-first score ordering. Later regression evidence preserves its
binary conclusions. Scores and AUC remain attached to their recorded source
version. Keeping both provenance boundaries is preferable to silently
rewriting old results as if all measurements came from the current
source tree. The generated manifest makes this distinction auditable.

\section{Deployment interpretation}
\label{sec:deployment-discussion}

The constrained Docker study supports a limited but useful statement:
one supplied-window process runs within the configured CPU and memory
limits on the measured Linux/amd64 host, including repeated rejected
updates. This is consistent with the motivation for nearby network
processing in edge systems \cite{shi2016}. ARM gateways, routers, and
microcontrollers require direct measurements of their instruction set,
memory hierarchy, storage, and power constraints.

The runtime receives completed counts rather than capturing packets
inside the measured service. Live collection could dominate CPU,
memory, and packet-loss behaviour. It must also preserve the exact
device filter, protocol exclusivity, window alignment, and missingness
semantics used to train the model. A fast collector with different
semantics would feed a different statistical problem to a valid
saved model.

State publication and bounded logging improve operational control.
Input authenticity, remote transport, downstream retention,
operator notification, and incident response remain integration
responsibilities. Durable delivery requires an acknowledgement from the
external consumer. These boundaries are material to deployment and define
work beyond the supplied-window experiment.

Finally, deployment should distinguish the quality of a decision
from the reliability of its execution. The implementation can
validate inputs, preserve state consistency, and reproduce a
decision while attacks with unchanged four-count representations remain
invisible. Conversely, a useful statistical signal is insufficient
if a collector drops windows or an alert is never delivered. The
thesis supplies evidence for both the statistical and software
components within stated limits, without treating either as a
substitute for the other.
